\section{Discussion}
\label{sec:discussion}

\paragraph{What the evidence establishes.}
On 15 small, fully numerical classification datasets, the \texttt{small} Kumo Tabular checkpoint with default settings
produced better probabilistic predictions than each of seven untuned classical models, consistently across datasets,
with differences that survive multiplicity correction and both sensitivity analyses (RQ1, RQ2). The advantage also
appears in ROC-AUC, so it is not explained by calibration alone. On six regression datasets it ranked first on average,
but the evidence is too limited for a statistical claim.

\paragraph{Where it does not help.}
The aggregate hides dataset-level exceptions (RQ3). Kumo Tabular's accuracy was below the best baseline on three of 15
classification datasets, and its RMSE was 44\% higher than XGBoost's on \texttt{energy\_efficiency}. In the case
studies, logistic regression matched it for classification with 25 to 100 training rows, and on the diabetes data
linear regression was statistically indistinguishable from it. The largest accuracy gains occurred on two multiclass
datasets (\texttt{vehicle}, \texttt{balance-scale}); on several binary tasks the differences in accuracy were a fraction
of a percentage point.

\paragraph{Interpreting the comparison.}
The baselines were not tuned, whereas Kumo Tabular's defaults result from extensive pretraining. Our comparison
therefore answers a practical question---what a practitioner obtains without model selection on small numerical
tables---and not which model class is better when each is tuned. Regularized or tuned gradient boosting, and a strong
linear model for nearly linear targets, remain relevant alternatives.

\paragraph{Cost.}
Prediction with Kumo Tabular processes the full context at every call. On our hardware its time per fold was comparable
to large random forests and smaller than CatBoost's, but well above linear models, XGBoost and LightGBM, and it requires
a GPU for reasonable speed (RQ4).

\section{Conclusion}
\label{sec:conclusion}
In an independent evaluation on 21 small numerical OpenML datasets, the smallest Kumo Tabular checkpoint
outperformed seven untuned classical baselines on classification log loss on nearly every dataset, and ranked first on
average for regression without statistical support. The results justify Kumo Tabular as a strong default for small
numerical tables, while the dataset-level exceptions, the restricted data regime and the untuned baselines limit how
far the conclusion extends.
